% Standalone note. From the repository root, compile with:
% latexmk -pdf -outdir=.vscode/latex-build/notes notes/minimax-regret-counterexample.tex
\documentclass[11pt]{article}
\usepackage[T1]{fontenc}
\usepackage{times}
\usepackage{amsmath,amssymb,amsthm}
\usepackage[margin=1in]{geometry}

\newtheorem{theorem}{Theorem}
\newcommand{\ind}{\mathbf 1}
\newcommand{\E}{\mathbb E}
\newcommand{\V}{\mathcal V}
\newcommand{\C}{\mathcal C}
\newcommand{\X}{\mathcal X}
\newcommand{\Lo}{\mathcal L}
\newcommand{\Reg}{\mathcal R}
\newcommand{\Mi}{\mathcal M}

\title{Worst-case regret optimality does not guarantee accurate predictions}
\author{}
\date{}

\begin{document}
\maketitle

We give a fixed verifier class with two certificates and two hypotheses.
There is a pair of debaters that each exactly minimizes its own worst-case
expected regret for every horizon, but makes $T/2$ expected prediction
mistakes. Cooperative learners make at most one mistake.
The example uses the freedom to choose among minimax-optimal strategies.
Its proposer deliberately makes unnecessary errors on easy inputs.
Thus the claim is that minimax optimality alone does not guarantee accurate
predictions, not that every minimax-optimal pair must fail.

\section{The verifier class}

Let $\X=\{e\}\cup\mathbb N$, where $e\notin\mathbb N$, and let
$\C=\{0,1\}$. Debates have length one. For $b\in\{0,1\}$, define
\[
h_b(x)=
\begin{cases}
b,&x=e,\\
0,&x\in\mathbb N.
\end{cases}
\]
For each $b\in\{0,1\}$ and $s:\mathbb N\to\{0,1\}$, define
\[
v_{b,s}(x,p,r)=
\begin{cases}
b,&x=e,\\
0,&x\in\mathbb N\text{ and }(p=s(x)\text{ or }r=0),\\
1,&\text{otherwise}.
\end{cases}
\]
Set $\V=\{v_{b,s}:b\in\{0,1\},\ s:\mathbb N\to\{0,1\}\}$.
At a natural-number input, the verifier has the table
\[
\begin{array}{c|cc}
 &r=0&r=1\\\hline
p=s(x)&0&0\\
p\ne s(x)&0&1
\end{array}
\]
The certificate $p=s(x)$ forces label $0$, and $r=0$ prevents any
certificate from forcing label $1$. At $e$, every entry is $b$.
Consequently $v_{b,s}$ validates exactly $h_b$.

\section{The objective and the separation}

Write $\hat y_t$ for the proposed label and $\tilde y_t$ for the verifier
feedback. As in the manuscript, define
\[
\Lo_P(T)=\sum_{t=1}^T\ind[\tilde y_t\ne\hat y_t],
\qquad
\Lo_R(T)=\sum_{t=1}^T\ind[\tilde y_t=\hat y_t].
\]
A round is irrefutable when $v_{b,s}(x_t,p_t,r)=\hat y_t$ for every
$r\in\C$. Let $U(T)$ count these rounds, and set
$\Reg_P(T)=\Lo_P(T)$ and $\Reg_R(T)=\Lo_R(T)-U(T)$.
Prediction mistakes are
$\Mi(T)=\sum_{t=1}^T\ind[\hat y_t\ne h_b(x_t)]$.

Each debater minimizes the supremum of its expected regret over fixed
verifiers in $\V$, adaptive input sequences, and strategies of the other
debater. Expectations are over private randomization. The environment
chooses $x_t$ before the proposer's current action, and the responder sees
that action before responding. In particular, this is the manuscript's
separate worst-case objective, rather than a common-prior objective or a
best response to one fixed opponent.

\begin{theorem}
For this fixed verifier class, there are proposer and responder strategies
that exactly minimize their respective worst-case expected regrets for
every horizon $T$. Their worst-case expected regrets are $T/2$ and $0$,
respectively. On the input sequence $x_t=e$ for all $t$, the pair makes
$T/2$ expected prediction mistakes. There is a cooperative learner with
at most one prediction mistake on every input sequence.
\end{theorem}
\begin{proof}
We first give a responder that refutes every refutable prediction. It
always chooses $r_t=1-\hat y_t$. At a natural-number input, a prediction
of $1$ is refuted by $r_t=0$. A prediction of $0$ is refuted by $r_t=1$
unless $p_t=s(x_t)$, in which case the round is irrefutable. At $e$, an
incorrect prediction is automatically refuted and a correct prediction is
irrefutable. Thus $\Reg_R(T)=0$ on every run, against every proposer.
Since responder regret is nonnegative, this responder is minimax-optimal.

The proposer uses fresh independent fair coins. On a natural-number
input, it predicts $0$ and chooses $p_t$ uniformly from $\{0,1\}$.
On input $e$, it predicts a uniformly random bit and chooses $p_t=0$.
At a natural-number input, the event $p_t=s(x_t)$ has probability $1/2$
and guarantees zero proposer loss, regardless of the responder's action.
At $e$, the feedback is always $b$, so the expected loss is $1/2$.
Hence this proposer guarantees $\E\Reg_P(T)\le T/2$ against every
verifier, adaptive environment, and responder.

To prove that this guarantee is optimal, fix any proposer. Present the
distinct inputs $x_t=t$, let the responder always choose $r_t=1$, and
sample $s(1),\ldots,s(T)$ independently and uniformly before play.
The value of $b$ can be fixed arbitrarily. Conditional on the history and
the proposer's current label and certificate, the fresh value $s(t)$ is
still a fair bit. Therefore the feedback
$\tilde y_t=\ind[p_t\ne s(t)]$ is a fair bit, and the proposer's expected
loss in each round is $1/2$. Its expected regret, averaged over these
fixed verifiers, is $T/2$. Some fixed verifier therefore gives expected
regret at least $T/2$. This proves the minimax lower bound.

Now pair the two specified strategies and present $e$ in every round.
The proposer predicts a fresh fair bit each time, so
$\E\Mi(T)=T/2$ for either fixed value of $b$. Both strategies remain
exactly minimax-optimal by the preceding arguments.

Finally, a cooperative learner predicts $0$ on every natural-number
input. On the first occurrence of $e$, it predicts arbitrarily and
observes feedback $b$. It predicts $b$ on every later occurrence of $e$.
Only the first occurrence of $e$ can be a prediction mistake.
\end{proof}

\section{What the example establishes}

The proposer learns $b$ on the first occurrence of $e$ but continues to
predict randomly there. Replacing that behavior with prediction $b$
improves its performance on these inputs without improving its worst-case
regret, because fresh natural-number inputs still force expected regret
$T/2$. Both behaviors are therefore minimax-optimal.

This is a separation between what cooperation can achieve and what
separate minimax optimality alone guarantees. It does not establish
failure for every choice among minimax-optimal strategies, for strategies
required to remain optimal conditional on each history, or for mutual
best responses on the realized input sequence. Nor are both debaters
no-regret learners in the sense of having sublinear regret. The
proposer's optimal worst-case regret is linear, consistently with the
manuscript's inequality $\Mi(T)\le\Reg_P(T)+\Reg_R(T)$.

\end{document}
